% 原书物理页=印刷页1--6；止于1.1.5之前.
\chapter{解析理论}\label{chap:modular-1}
\section{Poincaré上半平面及其商}\label{sec:modular-1-1}
令 $\mathbb H:=\{z\in\mathbb C,\Im(z)>0\}$. 这是 $\mathbb C$ 中连通且单连通的开集，由映射 $z\mapsto (z-i)/(z+i)$ 与开单位圆盘 $\mathbb D$ 双全纯同构. Riemann曲面理论表明，在双全纯同构意义下，单连通的\newterm{modular-riemann-surface}{Riemann曲面}{Riemann surface}只有 $\mathbb D\simeq\mathbb H$、复平面 $\mathbb C$ 和 Riemann球面 $\mathbb P^1(\mathbb C)=(\mathbb C^2\setminus\{0\})/\mathbb C^\times$.

\subsection{$\mathbb H$ 的自同构}\label{sec:modular-1-1-1}
$\mathrm{GL}_2(\mathbb C)$ 在 $\mathbb C^2$ 上的线性作用下降为 $\mathbb P^1(\mathbb C)$ 上的双全纯作用，中心 $\mathbb C^\times$ 的作用是平凡的. 写成 $\mathbb P^1(\mathbb C)=\mathbb C\coprod\{\infty\}$，其中借助 $z\mapsto(1,z)\mathbb C^\times$ 将 $\mathbb C$ 嵌入 $\mathbb P^1(\mathbb C)$.\footnote{译注：原文的坐标写作 $(1,z)$，而下文的分式线性作用按通常列向量约定对应 $(z,1)$；此处保留原文.}
这一坐标给出分解 $\mathbb P^1(\mathbb C)=\mathbb H\coprod\mathbb P^1(\mathbb R)\coprod\overline{\mathbb H}$，行列式为正的矩阵组成的子群 $\mathrm{GL}_2^+(\mathbb R)$ 保持这一分解，特别地作用于 $\mathbb H$. 具体说，若 $\gamma=\begin{bmatrix}a&b\\c&d\end{bmatrix}$，则
\begin{align*}
\gamma z=\frac{az+b}{cz+d},\qquad \Im(\gamma z)=\frac{\det(\gamma)}{|cz+d|^2}\Im(z).
\end{align*}
\begin{theorem*}
设 $\operatorname{Aut}(\mathbb H)$ 为 $\mathbb H$ 的双全纯自同构群.
\begin{enumerate}[label=\roman*)]
\item 上述作用诱导群同构
\begin{align*}
\mathrm{SL}_2(\mathbb R)/\{\pm1\}=\mathrm{GL}_2^+(\mathbb R)/\mathbb R^\times\xrightarrow{\sim}\operatorname{Aut}(\mathbb H).
\end{align*}
\item $\mathrm{SL}_2(\mathbb R)$ 在 $\mathbb H$ 上的作用是传递的，$i$ 在 $\mathrm{SL}_2(\mathbb R)$ 中的稳定子为 $\mathrm{SO}_2(\mathbb R)$. 此外，$\gamma\mapsto\gamma i$ 诱导同胚
\begin{align*}
\mathrm{SL}_2(\mathbb R)/\mathrm{SO}_2(\mathbb R)\xrightarrow{\sim}\mathbb H.
\end{align*}
\end{enumerate}
\end{theorem*}
在 ii) 中，$\mathrm{SL}_2(\mathbb R)/\mathrm{SO}_2(\mathbb R)$ 配备商拓扑，即使典范投影 $\mathrm{SL}_2(\mathbb R)\to\mathrm{SL}_2(\mathbb R)/\mathrm{SO}_2(\mathbb R)$ 连续的最细拓扑.
\begin{proof}
先证 ii). 计算可知 $\begin{bmatrix}a&b\\c&d\end{bmatrix}$ 固定 $i$ 当且仅当 $a=d$ 且 $b=-c$. 行列式条件又给出 $a^2+b^2=1$，故 $i$ 在 $\mathrm{SL}_2(\mathbb R)$ 中的稳定子为 $\mathrm{SO}_2(\mathbb R)$. 另一方面，写 $z=x+iy$，则
\begin{align*}
z=\begin{bmatrix}y&x\\0&1\end{bmatrix}i=\begin{bmatrix}y^{-1/2}&-xy^{-1/2}\\0&y^{1/2}\end{bmatrix}i,
\end{align*}
由此得到作用的传递性.\footnote{译注：原文第二个矩阵作用于 $i$ 所得是 $(i-x)/y$；给出 $x+iy$ 的行列式为 $1$ 的矩阵应为 $\left[\begin{smallmatrix}y^{1/2}&xy^{-1/2}\\0&y^{-1/2}\end{smallmatrix}\right]$. 原式照录.}
因此 ii) 中的映射是双射. $\mathrm{SL}_2(\mathbb R)$ 在 $\mathbb H$ 上的作用显然连续，故此双射连续，只需再证它是开映射.\footnote{当拓扑群 $G$ 在空间 $X$ 上连续且传递地作用时，对每个 $x\in X$，$X$ 上的作用诱导连续双射 $G/G_x\to X$. 它未必是同胚，但若 $G$ 与 $X$ 都局部紧，则它总是同胚.}
为此，只须证明若 $U$ 是 $\mathrm{SL}_2(\mathbb R)$ 的开集，则 $U.i$ 是 $i$ 在 $\mathbb H$ 中的邻域；由公式
\begin{align*}
x+iy=\begin{bmatrix}y^{1/2}&-xy^{-1/2}\\0&y^{-1/2}\end{bmatrix}i
\end{align*}
可见这一点.\footnote{译注：原文这里的负号应为正号；此外，关于 $U.i$ 为 $i$ 的邻域的说法须在 $U$ 含单位元时理解.}
再证 i). 计算可知作用的核确为 $\{\pm1\}$. 为证满射，取 $\gamma\in\operatorname{Aut}(\mathbb H)$. 由 ii)，必要时与某个 $\alpha\in\mathrm{SL}_2(\mathbb R)$ 复合后，可设 $\gamma i=i$. 通过
\begin{align*}
\tau=\begin{bmatrix}1&-i\\1&i\end{bmatrix}:z\longmapsto\frac{z-i}{z+i}
\end{align*}
转移这一作用，得到开单位圆盘 $\mathbb D$ 上将 $0$ 映到 $0$ 的自同构 $\tau\gamma\tau^{-1}$. Schwarz 引理表明它是 $z\mapsto\zeta z$ 的伸缩变换，其中 $|\zeta|=1$. 因此对某个 $\theta\in\mathbb R$，有
\begin{align*}
\tau\gamma\tau^{-1}(z)=\begin{bmatrix}e^{i\theta}&0\\0&e^{-i\theta}\end{bmatrix}z.
\end{align*}
计算得到
\begin{align*}
\tau^{-1}\begin{bmatrix}e^{i\theta}&0\\0&e^{-i\theta}\end{bmatrix}\tau=\begin{bmatrix}\cos\theta&-\sin\theta\\\sin\theta&\cos\theta\end{bmatrix},
\end{align*}
故 $\gamma$ 与右侧矩阵在 $\mathbb H$ 上诱导相同的自同构.
\end{proof}
$\gamma\in\mathrm{SL}_2(\mathbb R)$ 在 $\mathbb C$ 中有两个特征值（计重数）. 称 $\gamma$ 为
\begin{itemize}
\item \newterm{modular-elliptic-element}{椭圆型}{elliptic}，若其特征值均非实数；
\item \newterm{modular-hyperbolic-element}{双曲型}{hyperbolic}，若其特征值为两个不同的实数；
\item \newterm{modular-parabolic-element}{抛物型}{parabolic}，若其特征值相等（因而为实数），且 $\gamma\ne\pm1$.
\end{itemize}
考察特征多项式可得：$\gamma$ 为椭圆型当且仅当 $|\operatorname{tr}(\gamma)|<2$；为双曲型当且仅当 $|\operatorname{tr}(\gamma)|>2$；为抛物型当且仅当 $|\operatorname{tr}(\gamma)|=2$ 且 $\gamma\ne\pm1$. 因此每个 $\gamma\ne\pm1$ 恰属于这三类之一. 也可通过 $\gamma$ 在 $\mathbb P^1(\mathbb C)$ 中的不动点刻画这些类型.
\begin{lemma*}
设 $\gamma\in\mathrm{SL}_2(\mathbb R)$.
\begin{enumerate}[label=\roman*)]
\item $\gamma$ 为椭圆型当且仅当它在 $\mathbb H$ 中有不动点. 此不动点唯一，且 $\gamma$ 在 $\mathrm{SL}_2(\mathbb R)$ 中共轭于旋转 $\left[\begin{smallmatrix}\cos\theta&-\sin\theta\\\sin\theta&\cos\theta\end{smallmatrix}\right]\in\mathrm{SO}_2(\mathbb R)$.
\item $\gamma$ 为双曲型当且仅当它在 $\mathbb P^1(\mathbb R)$ 中有两个不动点.
\item $\gamma$ 为抛物型当且仅当它在 $\mathbb P^1(\mathbb R)$ 中有唯一不动点. 此外，$\gamma$ 在 $\mathrm{GL}_2(\mathbb R)$ 中共轭于形如 $\pm\left[\begin{smallmatrix}1&1\\0&1\end{smallmatrix}\right]$ 的矩阵.
\end{enumerate}
\end{lemma*}
\noindent\textbf{译注：}原文引理未排除 $\gamma=\pm1$；这两个元素固定所有点，故 i) 的逆向需要这一排除，ii) 中的``两个''也应理解为``恰有两个''.
\begin{proof}
i) $\gamma$ 为椭圆型当且仅当作为 $\mathrm{GL}_2(\mathbb R)$ 的元素，它在 $\mathbb C$ 上可对角化而在 $\mathbb R$ 上不可对角化. 此时它在 $\mathbb C^2$ 中恰有两条特征直线，在 $\mathbb R^2$ 中则没有. 因而它不固定 $\mathbb P^1(\mathbb R)$ 中的点，而在 $\mathbb P^1(\mathbb C)\setminus\mathbb P^1(\mathbb R)$ 中固定两个共轭点，其中恰有一个在 $\mathbb H$ 中. 这个点在 $\mathrm{SL}_2(\mathbb R)$ 作用下与 $i$ 等价，故 $\gamma$ 共轭于 $i$ 的稳定子 $\mathrm{SO}_2(\mathbb R)$ 中的元素.

ii) 显然. iii) 依定义，$\gamma$ 为抛物型当且仅当它在 $\mathrm{GL}_2(\mathbb R)$ 中共轭于不等于 $\pm1$ 的矩阵 $\left[\begin{smallmatrix}a&b\\0&a\end{smallmatrix}\right]$. 由 $\det\gamma=1$ 得 $a=\pm1$. 用 $\left[\begin{smallmatrix}1&0\\0&\pm b\end{smallmatrix}\right]$ 共轭，便得到 $\pm\left[\begin{smallmatrix}1&1\\0&1\end{smallmatrix}\right]$. 后者在 $\mathbb P^1(\mathbb C)$ 中有唯一不动点 $\infty$.
\end{proof}

\subsection{$\mathrm{SL}_2(\mathbb R)$ 的离散子群}\label{sec:modular-1-1-2}
若 $\mathrm{SL}_2(\mathbb R)$ 的子群 $\Gamma$ 对诱导拓扑是离散的，就称它为\newterm{modular-discrete-subgroup}{离散子群}{discrete subgroup}. 第一个例子是 $\mathrm{SL}_2(\mathbb Z)$，称为\newterm{modular-modular-group}{模群}{modular group}，也记作 $\Gamma(1)$. 更一般地，对 $N\in\mathbb N$，记
\begin{align*}
\Gamma(N)&:=\operatorname{Ker}\bigl(\mathrm{SL}_2(\mathbb Z)\to\mathrm{SL}_2(\mathbb Z/N\mathbb Z)\bigr)\\
&=\left\{\begin{bmatrix}a&b\\c&d\end{bmatrix}\equiv\begin{bmatrix}1&0\\0&1\end{bmatrix}\pmod N\right\}.
\end{align*}
这是 $\mathrm{SL}_2(\mathbb Z)$ 的有限指数子群，称为\newterm{modular-principal-congruence-subgroup}{级为 $N$ 的主同余子群}{principal congruence subgroup of level $N$}. 本课程中特别关注以下群：
\begin{align*}
\Gamma_1(N)&:=\left\{\begin{bmatrix}a&b\\c&d\end{bmatrix}\equiv\begin{bmatrix}1&*\\0&1\end{bmatrix}\pmod N\right\},\\
\Gamma_0(N)&:=\left\{\begin{bmatrix}a&b\\c&d\end{bmatrix}\equiv\begin{bmatrix}*&*\\0&*\end{bmatrix}\pmod N\right\}.
\end{align*}
有 $\Gamma(N)\subset\Gamma_1(N)\subset\Gamma_0(N)\subset\Gamma(1)=\mathrm{SL}_2(\mathbb Z)$.
$\mathrm{SL}_2(\mathbb R)$ 还有许多不与 $\mathrm{SL}_2(\mathbb Z)$ ``可公度''的离散子群：即 $\Gamma\cap\mathrm{SL}_2(\mathbb Z)$ 既不是 $\mathrm{SL}_2(\mathbb Z)$ 的有限指数子群，也不是 $\Gamma$ 的有限指数子群. 其中一些由四元数代数得到，在数论中也起着重要作用.
\begin{proposition*}\label{prop:modular-discrete-action}
设 $\Gamma$ 为 $\mathrm{SL}_2(\mathbb R)$ 的离散子群.
\begin{enumerate}[label=\roman*)]
\item 对任意 $z\in\mathbb H$，$z$ 在 $\Gamma$ 中的稳定子 $\Gamma_z$ 是循环群（特别地，是有限群）.
\item 每个 $z\in\mathbb H$ 都有开邻域 $U$，使得对所有 $\gamma\in\Gamma$，$\gamma U\cap U\ne\varnothing\Rightarrow\gamma z=z$.
\item 两条不同的轨道 $\Gamma z$ 和 $\Gamma z'$ 有互不相交的开邻域.
\end{enumerate}
\end{proposition*}
\begin{proof}
i) $z$ 的稳定子共轭于 $\mathrm{SO}_2(\mathbb R)$ 的离散子群，故为有限群. 这样的子群是循环群（将 $\mathrm{SO}_2(\mathbb R)$ 与 $\mathbb C^\times$ 中的单位圆周等同）.

为证 ii) 和 iii)，使用如下事实：若 $A,B$ 为 $\mathbb H$ 中的两个紧集，则
\begin{align*}
\mathrm{SL}_2(\mathbb R)_{A,B}:=\{\gamma\in\mathrm{SL}_2(\mathbb R),\ \gamma A\cap B\ne\varnothing\}
\end{align*}
是 $\mathrm{SL}_2(\mathbb R)$ 中的紧集. 事实上，若以 $\pi$ 表示映射 $\gamma\mapsto\gamma i$，以 $\iota$ 表示它的截面
\begin{align*}
x+iy\longmapsto\begin{bmatrix}y^{1/2}&-xy^{-1/2}\\0&y^{-1/2}\end{bmatrix},
\end{align*}
则
\begin{align*}
\mathrm{SL}_2(\mathbb R)_{A,B}=\{\gamma\in\mathrm{SL}_2(\mathbb R),\ \gamma\pi^{-1}(A)\cap\pi^{-1}(B)\ne\varnothing\}
\end{align*}
是紧的，因为 $\pi^{-1}(A)=\iota(A)\mathrm{SO}_2(\mathbb R)$ 和 $\pi(B)$ 都是紧的.\footnote{译注：原文截面矩阵的右上元仍带负号，应为正号；末尾 $\pi(B)$ 应为 $\pi^{-1}(B)$. 此处均保留原式.}

ii) 将此用于 $z$ 的紧邻域 $V$. 集合 $\Gamma_V:=\{\gamma\in\Gamma,\ \gamma V\cap V\ne\varnothing\}$ 是 $\Gamma$ 中的紧集，故为有限集. 对每个不固定 $z$ 的 $\gamma\in\Gamma_V$，可取 $z$ 的开邻域 $U_\gamma$，使得 $\gamma U_\gamma\cap U_\gamma=\varnothing$. 对这样的 $\gamma$ 取有限交，便得到所需的 $U$.

iii) 同上，若 $V'$ 是 $z'$ 的紧邻域，则集合 $\Gamma_{V,V'}:=\{\gamma\in\Gamma,\ \gamma V\cap V'\ne\varnothing\}$ 有限，由此可取 $z$ 的开邻域 $U$ 和 $z'$ 的开邻域 $U'$，使得对所有 $\gamma$ 都有 $\gamma U\cap U'=\varnothing$. 于是 $\bigcup_\gamma\gamma U$ 和 $\bigcup_\gamma\gamma U'$ 分别是 $\Gamma z$ 和 $\Gamma z'$ 的两个不交邻域.
\end{proof}
命题的 iii) 说明商拓扑空间 $\Gamma\backslash\mathbb H$ 是 Hausdorff 的. 我们希望赋予它复结构. 当 $\Gamma$（更确切地说，$\Gamma/(\Gamma\cap\{\pm1\})$）在 $\mathbb H$ 上自由作用时（即对所有 $\gamma\in\Gamma$ 和 $z\in\mathbb H$，$\gamma z=z\Rightarrow\gamma=1$），这是容易且经典的. 然而 $\mathrm{PSL}_2(\mathbb Z)$ 的作用并不自由！\footnote{译注：这里及下文关于自由作用的说法应按有效商群 $\Gamma/(\Gamma\cap\{\pm1\})$ 理解；原文括号内仍写 $\gamma\in\Gamma$，照录.}

\subsection{作为Riemann曲面的 $\Gamma\backslash\mathbb H$}\label{sec:modular-1-1-3}
对 $z\in\mathbb H$，记 $h_z$ 为 $\Gamma_z/(\Gamma_z\cap\{\pm1\})$ 的阶. 若 $h_z>1$，则称 $z$ 为 $\Gamma$ 的\newterm{modular-elliptic-point}{椭圆点}{elliptic point}. 因此，$\Gamma$ 自由作用当且仅当它没有椭圆点.
回顾一下，要在拓扑空间 $X$ 上定义一维复结构，可以给定一个全纯图册（或它的等价类），也可以给定函数层 $\mathcal O_X$，使得每个 $x\in X$ 都有开邻域 $U$，且带环空间 $(U,(\mathcal O_X)|_U)$ 同构于 $\mathbb C$ 中配备全纯函数层的某个开集.
\begin{proposition*}
记 $\pi:\mathbb H\to\Gamma\backslash\mathbb H$ 为典范投影，对 $\Gamma\backslash\mathbb H$ 的每个开集 $U$，令
\begin{align*}
\mathcal O_{\Gamma\backslash\mathbb H}(U):=\mathcal O_{\mathbb H}(\pi^{-1}(U))^\Gamma
\end{align*}
（即 $\pi^{-1}(U)$ 上在 $\Gamma$ 作用下不变的全纯函数）. 则 $\mathcal O_{\Gamma\backslash\mathbb H}$ 是一个层，并在 $\Gamma\backslash\mathbb H$ 上定义使 $\pi$ 全纯的复结构. 这一Riemann曲面也记作 $Y(\Gamma)$.
\end{proposition*}
\begin{proof}
$\mathcal O_{\Gamma\backslash\mathbb H}$ 显然是一个层. 若 $z$ 不是椭圆点，可取 $z$ 的开邻域 $V$，使 $\pi|_V$ 是同胚 $V\xrightarrow{\sim}\pi(V)$. 此时 $\pi$ 诱导所需的同构 $(V,\mathcal O_V)\xrightarrow{\sim}(U,\mathcal O_U)$. 若 $z$ 是椭圆点，取 $\tau\in\mathrm{SL}_2(\mathbb C)$，使得 $\tau(\mathbb H)=\mathbb D$ 且 $\tau z=0$. Schwarz 引理保证 $\tau\Gamma_z\tau^{-1}$ 在 $\mathbb D$ 上通过旋转作用. 更确切地说，存在群同构 $\iota:\Gamma_z\xrightarrow{\sim}\mu_{h_z}$（$h_z$ 次单位根群），使得对所有 $z'\in\mathbb H$，有 $\tau(\gamma z')=\iota(\gamma)\tau(z')$.\footnote{译注：原文此同构的定义域应为有效商群 $\Gamma_z/(\Gamma_z\cap\{\pm1\})$；保留原式.}
对任意半径 $r<1$，开集 $V:=\tau^{-1}(D(0,r))$ 在 $\Gamma_z$ 作用下稳定. 由上一命题的 ii)，当 $r$ 充分小时，$\gamma V\cap V\ne\varnothing\Rightarrow\gamma\in\Gamma_z$. 于是有交换图
\[
\begin{tikzcd}[column sep=large,row sep=large]
V \arrow[r,"\sim","\tau"'] \arrow[d,"\pi"'] & D(0,r) \arrow[d,"x\mapsto x^{h_z}"]\\
\pi(V)=\Gamma_z\backslash V \arrow[r,dashed,"\sim","\bar\tau"'] & \mu_{h_z}\backslash D(0,r^{h_z})=D(0,r)
\end{tikzcd}
\]
其中 $\bar\tau$ 是同胚.\footnote{译注：图右下角按原图照录；与右侧 $x\mapsto x^{h_z}$ 及下文相容的目标应为 $\mu_{h_z}\backslash D(0,r)\simeq D(0,r^{h_z})$.}
由于对每个 $U\subset\pi(V)$ 都有 $\mathcal O(U)=\mathcal O(\pi^{-1}(U)\cap V)^{\Gamma_z}$，可见 $\bar\tau$ 将函数层 $\mathcal O_{\pi(V)}$ 转移为 $x\in D(0,r)$ 上在 $\mu_{h_z}$ 作用下不变的全纯函数层；它也就是关于 $x^{h_z}$ 的全纯函数层，因而是 $D(0,r^{h_z})$ 上通常的全纯函数层.
所以命题中的层确实在 $\Gamma\backslash\mathbb H$ 上定义复结构，而 $\pi$ 显然全纯. 特别要记住，$z'\mapsto\tau(z')^{h_z}$ 下降为 $\Gamma\backslash\mathbb H$ 中 $\pi(z)$ 附近的局部坐标.
\end{proof}

\subsection{尖点}\label{sec:modular-1-1-4}
若 $\mathbb P^1(\mathbb R)$ 中的点 $x$ 被 $\Gamma$ 的某个抛物型元素固定，则称它为 $\Gamma$ 的\newterm{modular-cusp}{尖点}{cusp}. 此时，取 $\tau\in\mathrm{SL}_2(\mathbb R)$ 使 $\tau\infty=x$，则 $\tau^{-1}\Gamma_x\tau$ 是 $\infty$ 的稳定子的离散子群，由此可知存在 $h>0$，使得
\begin{align*}
\tau^{-1}\Gamma_x\tau\cdot\{\pm1\}=\left\{\pm\begin{bmatrix}1&h\\0&1\end{bmatrix}^{m},\ m\in\mathbb Z\right\}.
\end{align*}
以 $P_\Gamma\subset\mathbb P^1(\mathbb R)$ 表示 $\Gamma$ 的尖点集，令 $\mathbb H^*=\mathbb H^*_\Gamma:=\mathbb H\sqcup P_\Gamma$. 在 $\mathbb H^*$ 上赋予由 $\mathbb H$ 的开集及形如 $V_{\tau,r}:=\tau(U_r\sqcup\{\infty\})$ 的集合生成的拓扑，其中
\begin{itemize}
\item $r>0$，$U_r=\{z\in\mathbb H,\Im(z)>r\}$；
\item $\tau\in\mathrm{SL}_2(\mathbb R)$，且 $\tau\infty\in P_\Gamma$.
\end{itemize}
若写 $\tau=\begin{bmatrix}a&b\\c&d\end{bmatrix}$，则 $\tau(U_r)$ 是 $\mathbb H$ 中直径为 $r^{-1}c^{-2}$ 的开圆盘，其在 $\mathbb P^1(\mathbb C)$ 中的闭包含有 $\tau\infty=a/b$.\footnote{译注：原文 $a/b$ 应为 $a/c$；圆盘描述限于 $c\ne0$，$c=0$ 时像仍为水平上半平面.}
由构造，$\Gamma$ 在 $\mathbb H^*$ 上的作用保持此拓扑（即作用连续），故可在商 $\Gamma\backslash\mathbb H^*$ 上赋予商拓扑.
\begin{lemma*}
若 $x_1,x_2\in P_\Gamma$，则存在 $x_1$ 的邻域 $V_{\tau_1,r}$ 和 $x_2$ 的邻域 $V_{\tau_2,r}$，使得
\begin{align*}
\forall\gamma\in\Gamma,\quad \gamma V_{\tau_1,r}\cap V_{\tau_2,r}\ne\varnothing\ \Longrightarrow\ \gamma x_1=x_2.
\end{align*}
特别地，空间 $\Gamma\backslash\mathbb H^*$ 是 Hausdorff 的，子集 $\Gamma\backslash P_\Gamma$ 在其中是离散的.
\end{lemma*}
\begin{proof}
必要时将整个情形作共轭变换，可以设 $x_2=\infty$，从而 $V_{\tau_2,r}=U_r$.\footnote{译注：按定义 $V_{\tau_2,r}$ 还含尖点 $\infty$，原文此处略去它；以下关于交集的计算针对上半平面部分.}
此时有 $h_\infty>0$，使 $\Gamma_\infty\cdot\{\pm1\}=\{\pm\gamma_\infty^m,\ m\in\mathbb Z\}$，其中 $\gamma_\infty=\left[\begin{smallmatrix}1&h_\infty\\0&1\end{smallmatrix}\right]$. 记 $x=x_1$，并固定 $\tau=\tau_1$ 使 $\tau\infty=x$，便有前述 $h>0$. 我们将验证每个满足 $r^2>hh_\infty$ 的 $r$ 都适用.
为此，设 $\gamma\in\Gamma$，且 $\gamma\tau U_r\cap U_r\ne\varnothing$. 需要证明 $\gamma x=\gamma\tau\infty=\infty$. 以 $c(\sigma)$ 表示矩阵 $\sigma\in\mathrm{SL}_2(\mathbb R)$ 的 $(2,1)$ 元，只须证明 $c(\gamma\tau)=0$. 由 $\gamma\tau U_r$ 的描述，交集非空蕴含 $r<r^{-1}c(\gamma\tau)^{-2}$，即 $c(\gamma\tau)^2<r^{-2}$.
令 $\delta_1:=\tau^{-1}\gamma^{-1}\gamma_\infty\gamma\tau$，这是离散群 $\tau^{-1}\Gamma\tau$ 的元素. 计算得到 $c(\delta_1)=-c(\gamma\tau)^2h_\infty$，由 $r$ 的选取知 $|c(\delta_1)h|<1$. 需要证明 $c(\delta_1)=0$，这等价于 $\delta_1\in(\tau^{-1}\Gamma\tau)_\infty=\tau^{-1}\Gamma_x\tau$.
为此，令 $\delta_{n+1}:=\delta_n\left[\begin{smallmatrix}1&h\\0&1\end{smallmatrix}\right]\delta_n^{-1}$，它属于 $\tau^{-1}\Gamma\tau$. 同样的计算给出
\begin{align*}
c(\delta_{n+1})&=-c(\delta_n)^2h,\\
b(\delta_{n+1})&=a(\delta_n)^2h,\\
a(\delta_{n+1})&=1-a(\delta_n)c(\delta_n)h.
\end{align*}
因此 $c(\delta_n)=-c(\delta_1)(c(\delta_1)h)^{2^n-1}$ 当 $n\to\infty$ 时趋于 $0$，进而 $\delta_n$ 趋于矩阵 $\left[\begin{smallmatrix}1&h\\0&1\end{smallmatrix}\right]$.\footnote{译注：原文迭代公式照录；由前一递推式，$n\ge2$ 时应为 $c(\delta_n)=-c(\delta_1)(c(\delta_1)h)^{2^{n-1}-1}$. 另外，若 $\Gamma$ 不含 $-1$，上述抛物矩阵可能只在 $\tau^{-1}\Gamma\tau\cdot\{\pm1\}$ 中，证明的群归属陈述须在此有效作用意义下理解.}
由于 $\tau^{-1}\Gamma\tau$ 离散，$n$ 充分大时 $\delta_n=\left[\begin{smallmatrix}1&h\\0&1\end{smallmatrix}\right]$，故 $c(\delta_n)=0$，最终得到 $c(\delta_1)=0$. 所述推论留给读者证明.
\end{proof}
现在考虑映射 $z\mapsto\exp(2i\pi z/h)$. 它给出将 $\infty$ 映到 $0$ 的同胚
\begin{align*}
\begin{bmatrix}1&h\\0&1\end{bmatrix}^{\mathbb Z}\backslash(U_r\sqcup\{\infty\})\xrightarrow{\sim}D(0,\rho),\qquad\rho=\exp(-2\pi r/h),
\end{align*}
从而赋予 $U_r\sqcup\{\infty\}$ 复结构.\footnote{译注：原文这里应指上述商空间；$U_r\sqcup\{\infty\}$ 本身不能由该同胚直接取得所称复结构.}
若 $x=\tau\infty$ 是 $\Gamma$ 的尖点，且 $h$ 如本分目开头所述，则上一引理表明，当 $r$ 充分大时，有图
\[
\begin{tikzcd}[column sep=large,row sep=large]
V_{\tau,r} \arrow[r,"\sim","\tau^{-1}"'] \arrow[d,"\pi"'] & U_r\sqcup\{\infty\} \arrow[d,"z\mapsto\exp(2i\pi z/h)"]\\
\pi(V_{\tau,r})=\Gamma_x\backslash V_{\tau,r} \arrow[r,dashed,"\sim","\bar\tau^{-1}"'] & D(0,\rho)
\end{tikzcd}
\]
这就赋予 $\pi(V_{\tau,r})$ 复结构. 具体而言，函数 $\exp(2i\pi\tau(z)/h)$ 给出 $\pi(x)$ 附近的局部坐标.\footnote{译注：按本段约定 $x=\tau\infty$，局部坐标中的 $\tau(z)$ 应为 $\tau^{-1}(z)$；原式照录.}
这些结构与 $\Gamma\backslash\mathbb H$ 上的复结构可拼合，在 $\Gamma\backslash\mathbb H^*$ 上给出Riemann曲面结构；验证留给读者. 通常以 $X(\Gamma)$ 表示这一Riemann曲面.


\subsection{模商}\label{sec:modular-1-1-5}

我们来具体描述模群 $\Gamma(1)$ 的商 $\Gamma\backslash\mathbb H^*$.依照惯例，简记 $X(1):=X(\Gamma(1))$，$Y(1):=Y(\Gamma(1))$.

\begin{definition*}

$\Gamma$ 在 $\mathbb H$ 中的\newterm{modular-fundamental-domain}{基本域}{fundamental domain}是满足下列条件的连通开集 $D$：

\begin{itemize}

\item 对任意 $\gamma\in\Gamma$，$\gamma D\cap D\ne\varnothing$ 蕴含 $\gamma=\pm1$；

\item $\mathbb H=\bigcup_{\gamma\in\Gamma}\gamma\overline D$，其中 $\overline D$ 表示 $D$ 的闭包.

\end{itemize}

\end{definition*}

$\mathrm{SL}_2(\mathbb R)$ 的每个离散子群都有基本域.这源于 $\mathbb H$ 上存在 $\mathrm{SL}_2(\mathbb R)$ 不变的距离 $d(z,z')$.事实上，选定 $z_0$ 后，可以取\footnote{译注：此处原文未限制 $z_0$；要使所给集合成为基本域，须选取模去 $\{\pm1\}$ 后稳定子平凡的 $z_0$.}
\begin{align*}
D=\{z\in\mathbb H:\ \forall\gamma\in\Gamma\setminus\{\pm1\},\ d(z,z_0)<d(z,\gamma z_0)\}.
\end{align*}
距离 $d(z,z')$ 是 Poincaré 度量 $\frac{dz\,d\overline z}{|\operatorname{Im}(z)|^2}$ 的测地距离.容易看出，这个度量是 $\mathrm{SL}_2(\mathbb R)$ 不变的.对于该度量，测地线是垂直于实轴的半圆和竖直直线.

为了具体给出 $\Gamma(1)$ 的一个基本域，记
\begin{align*}
S:=\begin{pmatrix}0&-1\\1&0\end{pmatrix},\qquad T:=\begin{pmatrix}1&1\\0&1\end{pmatrix}.
\end{align*}
于是对每个 $z\in\mathbb H$，有 $Sz=-1/z$，$Tz=1+z$.还可看出，在 $\mathrm{PSL}_2(\mathbb R)=\Gamma(1)/\{\pm1\}$ 中有关系 $S^2=1$ 和 $(ST)^3=1$.\footnote{译注：原文此处写作 $\mathrm{PSL}_2(\mathbb R)$，应与 $\mathrm{PSL}_2(\mathbb Z)=\Gamma(1)/\{\pm1\}$ 区分；上述两个关系在两群中均成立.}

\begin{theorem*}

开集 $D:=\{z\in\mathbb H:|z|>1,\ |\operatorname{Re}(z)|<\frac12\}$ 是 $\Gamma(1)$ 的基本域.此外：

\begin{enumerate}[label=\roman*)]

\item 若 $z\ne z'\in\overline D$ 在 $\Gamma(1)$ 作用下共轭，则
\begin{align*}
&\operatorname{Re}z=\pm\tfrac12,\quad z'=z\pm1=T^{\pm1}z,\\
&\text{或}\quad |z|=1,\quad z'=-1/z=Sz.
\end{align*}
\item 元素 $z\in\overline D$ 的稳定子 $\Gamma(1)_z$ 严格包含 $\{\pm1\}$，当且仅当出现以下三种情形之一：
\begin{align*}
z&=i:\quad \Gamma(1)_i=\langle S\rangle,\\
&\hspace{1em}\Gamma(1)_i/\{\pm1\}\simeq\mathbb Z/2\mathbb Z;\\
z&=j:=\exp(2i\pi/3):\quad \Gamma(1)_j=\langle ST\rangle,\\
&\hspace{1em}\Gamma(1)_j/\{\pm1\}\simeq\mathbb Z/3\mathbb Z;\\
z&=-1/j=\exp(2i\pi/6):\quad \Gamma(1)_{j^2}=\langle TS\rangle,\\
&\hspace{1em}\Gamma(1)_{-1/j}/\{\pm1\}\simeq\mathbb Z/3\mathbb Z.
\end{align*}
\end{enumerate}

\end{theorem*}

\noindent\textit{译注：}第三行稳定子的下标 $j^2$ 沿用原文；与该行情形相应的下标应为 $-1/j$.

\begin{proof}

设 $\Gamma'$ 是 $\Gamma$ 中由 $S,T$ 生成的子群.先证明 $\mathbb H=\bigcup_{\gamma\in\Gamma'}\gamma D$.\footnote{译注：此处原文未写闭包；由以下论证实际得到的是 $\mathbb H=\bigcup_{\gamma\in\Gamma'}\gamma\overline D$.}对 $z\in\mathbb H$ 和 $\gamma=\begin{psmallmatrix}a&b\\c&d\end{psmallmatrix}$，公式 $\operatorname{Im}(\gamma z)=\operatorname{Im}(z)/|cz+d|^2$ 以及 $c,d$ 为整数这一事实表明，$\gamma\mapsto\operatorname{Im}(\gamma z)$ 能取到最大值.若最大值在 $z'=\gamma_0z$ 处取到，则 $|z'|\ge1$；否则就有 $\operatorname{Im}(Sz')=\operatorname{Im}(-1/z')>\operatorname{Im}(z')$.又因 $\operatorname{Im}(Tz')=\operatorname{Im}(z')$，必要时施以 $T$ 的作用，可以设 $|\operatorname{Re}(z')|\le1/2$.



现在证明 i) 和 ii).可以设 $\operatorname{Im}(z')\ge\operatorname{Im}(z)$.若 $z'=\gamma z$，则 $|cz+d|\le1$，故 $c=0,1$ 或 $-1$.必要时将 $\gamma$ 换成作用相同的 $-\gamma$，只需考虑 $c=0$ 和 $c=1$.若 $c=0$，则 $|d|=1$ 且 $\gamma z=T^bz$；这只可能在 $b=1,\operatorname{Re}(z)=-1/2$ 或 $b=-1,\operatorname{Re}(z)=1/2$ 时发生.\footnote{译注：原文此处未列出恒等作用 $b=0$；该情形在讨论 ii) 的稳定子时也须计入.}

若 $c=1$，因要求 $|z+d|\le1$，只有两种可能：$d=0$，或 $|d|=1$ 且 $z\in\{j,Sj\}$.若 $d=0$，由 $\det\gamma=1$ 得 $b=-1$，从而 $z'=a-1/z$.由于 $|z|=1$，$-1/z$ 属于 $\overline D$，所以同前有三种情形：$a=0$ 且 $|z|=1$；$a=1$ 且 $\operatorname{Re}(z)=-1/2$（此时 $z=j$）；$a=-1$ 且 $\operatorname{Re}(z)=1/2$（此时 $z=-1/j$）.若 $d=1,z=j$，则 $a-b=1$，故 $z'=a-1/(1+j)=a+j$，所以 $a=0$ 或 $1$.$d=-1,z=-1/j$ 的情形同理.

\end{proof}

\begin{remark*}

从几何上看，$D\cup\{\infty\}$ 是 $\mathbb H^*$ 中顶点为 $j,-1/j,\infty$ 的测地三角形.若考虑 $S\cdot D$，也许更清楚：它是顶点为 $j,-1/j,0$ 的测地三角形.

\end{remark*}

\begin{corollary*}

$\mathrm{PSL}_2(\mathbb Z)$ 由 $S,T$ 的像生成.

\end{corollary*}

\begin{proof}

同上一证明，记 $\Gamma'$ 为 $S,T$ 生成的群.已经看到 $\mathbb H=\bigcup_{\gamma'\in\Gamma'}\gamma'\overline D$.给定 $\gamma\in\Gamma$，选取 $z_0\in D$.可以找到 $\gamma'\in\Gamma'$ 使 $\gamma'\gamma z_0\in\overline D$.由 i) 得 $z_0=\gamma'\gamma z_0$，再由 ii) 得 $\gamma'\gamma=\pm1$.故 $\Gamma(1)/\{\pm1\}$ 由 $S,T$ 的像生成.

\end{proof}

\begin{remark*}

可以证明 $\langle S,T\mid S^2,(ST)^3\rangle$ 是 $\mathrm{PSL}_2(\mathbb Z)$ 的一个表示.特别地，$\mathrm{PSL}_2(\mathbb Z)$ 同构于 $\mathbb Z/2\mathbb Z$ 与 $\mathbb Z/3\mathbb Z$ 的自由积.

\end{remark*}

上述定理表明，$\Gamma(1)$ 有两个椭圆点轨道，一个的阶为 $2$，另一个的阶为 $3$：
\begin{align*}
\{\Gamma(1)\text{ 在 }\mathbb H\text{ 中的椭圆点}\}=\Gamma(1)i\sqcup\Gamma(1)j.
\end{align*}
至于尖点，显然 $\infty\in P_{\Gamma(1)}$，且
\begin{align*}
\Gamma(1)_\infty=\{\pm1\}\cdot\begin{pmatrix}1&1\\0&1\end{pmatrix}^{\mathbb Z}.
\end{align*}
事实上，$\mathrm{SL}_2(\mathbb Z)$ 的每个抛物元素都在 $\mathrm{GL}_2(\mathbb Q)$ 中共轭于 $\begin{psmallmatrix}1&1\\0&1\end{psmallmatrix}$.\footnote{译注：此断言沿用原文；迹为 $-2$ 的抛物元素应改为考虑其相反数，因共轭保持迹.两者在射影直线上的作用相同，故不影响以下尖点结论.}

因此，$\Gamma(1)$ 的每个尖点都属于 $\mathrm{GL}_2(\mathbb Q)\cdot\infty=\mathbb P^1(\mathbb Q)$.注意到（留作习题）$\mathbb P^1(\mathbb Q)=\mathrm{SL}_2(\mathbb Z)\infty$，便得
\begin{align*}
P_{\Gamma(1)}=\mathbb P^1(\mathbb Q)=\Gamma(1)\cdot\infty.
\end{align*}
\subsection{模商的 Riemann 曲面}\label{sec:modular-1-1-6}

\begin{corollary*}

Riemann 曲面 $Y(1)=\Gamma(1)\backslash\mathbb H$ 同构于复平面，而 $X(1)=\Gamma(1)\backslash\mathbb H$ 同构于 Riemann 球面.

\end{corollary*}

\noindent\textit{译注：}原文在 $X(1)$ 的商表达式中漏写星号；按定义应为 $\Gamma(1)\backslash\mathbb H^*$.

\begin{proof}

我们利用紧 Riemann 曲面的分类给出拓扑解释.以后将看到实现上述同构的具体函数.



只需证明 $X(1)$ 紧且亏格为 $0$.对紧性，由尖点的描述，注意到 $X(1)$ 是 $\overline D^*:=\overline D\sqcup\{\infty\}$ 在 $\pi:\mathbb H^*\longrightarrow\Gamma(1)\backslash\mathbb H^*$ 下的像.$\overline D^*$ 显然紧：给定一个开覆盖，其中一个开集形如 $(U_r\sqcup\{\infty\})\cap\overline D^*$，其补集在 $\overline D$ 中紧.



关于亏格为 $0$，有几种几何论证，这里不作形式化.下面给出其中一种.从拓扑上说，从 $\overline D$ 出发，经两步可得到 $Y(1)$：

\begin{itemize}

\item 将两条半直线 $\operatorname{Re}(z)=1/2$ 与 $\operatorname{Re}(z)=-1/2$ 等同；

\item 将``线段'' $[i,j]$ 与线段 $[i,Sj]$ 等同.

\end{itemize}

第一步给出一个圆柱，其底部的``圆''由测地线段 $[j,Sj]$ 的两端 $j,Sj$ 等同得到.第二步把这个圆``压平''为线段 $[i,j]$，从而封闭圆柱的一端.于是可以``看出''，$\overline D$ 到线段 $[j,Sj]$ 的``竖直''测地收缩下降到商，给出 $Y(1)$ 到 $[j,Sj]$ 的像的收缩.但该像是一条线段.因此所得空间单连通.于是 $Y(1)$ 同构于 $\mathbb C$ 或 $\mathbb D$；但因它可以通过添加一点紧化，所以只能是 $\mathbb C$.由此 $X(1)$ 同构于 $\mathbb P^1(\mathbb C)$.

\end{proof}



\subsection{同余子群}\label{sec:modular-1-1-7}

设 $\Gamma$ 为 $\Gamma(1)$ 的有限指数子群.



\paragraph{基本域}令 $d:=|\Gamma\{\pm1\}\backslash\Gamma(1)|$.选取 $\gamma_1,\ldots,\gamma_d$ 使 $\overline\Gamma(1)=\bigsqcup_i\overline\Gamma\,\overline\gamma_i$（横线表示在 $\mathrm{PSL}_2$ 中的像），则并集 $D'=\gamma_1D\sqcup\cdots\sqcup\gamma_dD$ 具有基本域的所有性质，唯独不连通.在具体情形中，关键就是选取 $\gamma_i$ 使闭包 $\overline D'$ 连通.此时，$\overline D'$ 的内部就是 $\Gamma$ 的一个基本域.

\paragraph{尖点}显然 $P_\Gamma\subset P_{\Gamma(1)}=\mathbb P^1(\mathbb Q)$.但每个尖点 $x\in\mathbb P^1(\mathbb Q)$ 的稳定子 $\Gamma(1)_x$ 都是无限群，而 $\Gamma$ 的指数有限，故 $\Gamma_x$ 也是无限群，$x$ 因而也是尖点.因此 $P_\Gamma=\mathbb P^1(\mathbb Q)$.不过，$\mathbb P^1(\mathbb Q)$ 中一般有多个 $\Gamma$ 轨道.记这些轨道的个数为 $n_\infty$.

\paragraph{椭圆点}若 $z$ 是 $\Gamma$ 的椭圆点，则它在 $\Gamma(1)$ 作用下共轭于 $i$ 或 $j$.记共轭于 $i$ 的椭圆点的 $\Gamma$ 轨道数为 $n_2$（这些点的阶 $h=2$），共轭于 $i$ 的椭圆点的 $\Gamma$ 轨道数为 $n_3$（这些点的阶 $h=3$）.\footnote{译注：定义 $n_3$ 时原文再次写 $i$；此处应为 $j$.}

\begin{theorem*}

Riemann 曲面 $X(\Gamma)$ 是紧的.投影 $X(\Gamma)\longrightarrow X(1)$ 是次数（或重数）为 $d$ 的全纯映射，且 $X(\Gamma)$ 的亏格为
\begin{align*}
g=1+d/12-n_2/4-n_3/3-n_\infty/2.
\end{align*}
\end{theorem*}

\begin{proof}

同 $\Gamma(1)$ 的情形，注意到 $X(\Gamma)$ 是紧集 $\gamma_1\overline D^*\cup\cdots\cup\gamma_d\overline D^*$ 的像，即得紧性.映射 $Y(N)\longrightarrow Y(1)$ 的全纯性源于复结构的定义，同理可得 $X(N)\longrightarrow X(1)$ 在尖点附近的全纯性.\footnote{译注：此处原文使用 $Y(N),X(N)$；本段所讨论的一般子群对应的记号应为 $Y(\Gamma),X(\Gamma)$.}

该映射非常值，故有一个次数（亦称重数）；可通过数一个``普通点''（既非椭圆点也非尖点）的原像来计算.由基本域的形状得到次数为 $d$.为计算亏格，使用 Hurwitz 公式
\begin{align*}
\chi(X(\Gamma))=d\cdot\chi(X(1))+\sum_P(e_P-1),
\end{align*}
其中 $\chi(X)=2g(X)-2$ 是紧 Riemann 曲面 $X$ 的 Euler 示性数，$e_P$ 表示 $f:X(\Gamma)\longrightarrow X(1)$ 在 $P\in X(\Gamma)$ 处的分歧指数（即在 $P$ 附近，$f$ 局部形如 $z\mapsto z^{e_P}$）.\footnote{译注：原文称 $2g-2$ 为 Euler 示性数；它实际是通常 Euler 示性数 $2-2g$ 的相反数.此处的 Hurwitz 公式与原文采用的符号约定一致.}

本情形由此得到
\begin{align*}
g(X(\Gamma))=1-d+\tfrac12\sum_P(e_P-1).
\end{align*}
记 $\pi:\mathbb H^*\longrightarrow X(1)$ 为典范投影.已经看到，$\pi$ 在 $\mathbb H$ 中一点处的分歧指数有限：普通点处为 $1$，等价于 $i$ 的椭圆点处为 $2$，等价于 $j$ 的椭圆点处为 $3$.由分歧指数的乘法性，$f$ 在一个非尖点 $P$ 处的分歧指数 $e_P$ 也只能是 $1,2$ 或 $3$.



具体而言，考虑纤维 $f^{-1}(\pi(i))$.对该纤维中的 $P$，有：

\begin{itemize}

\item 若 $P$ 是 $\Gamma$ 的椭圆点（的像），则 $e_P=1$；

\item 若 $P$ 是 $\Gamma$ 的普通点（的像），则 $e_P=2$.

\end{itemize}

由 $n_2$ 的定义，$f^{-1}(\pi(i))$ 中恰有 $n_2$ 个点满足 $e_P=1$.由次数的性质，该纤维中满足 $e_P=2$ 的点数因而为 $(d-n_2)/2$.所以它对 Hurwitz 公式的贡献为 $\sum_{P\mapsto\pi(i)}(e_P-1)=(d-n_2)/2$.



现在考虑纤维 $f^{-1}(\pi(j))$.对该纤维中的 $P$，有：

\begin{itemize}

\item 若 $P$ 是 $\Gamma$ 的椭圆点（的像），则 $e_P=1$；

\item 若 $P$ 是 $\Gamma$ 的普通点（的像），则 $e_P=3$.

\end{itemize}

由 $n_3$ 的定义，$f^{-1}(\pi(i))$ 中恰有 $n_3$ 个点满足 $e_P=1$.由次数的性质，$f^{-1}(\pi(i))$ 中满足 $e_P=3$ 的点数因而为 $(d-n_3)/3$.所以它对 Hurwitz 公式的贡献为 $\sum_{P\mapsto\pi(j)}(e_P-1)=2(d-n_2)/3$.\footnote{译注：本段两处 $\pi(i)$ 及末式的 $n_2$ 均沿用原文；对应当前纤维应分别为 $\pi(j)$ 及 $n_3$.}



\begingroup\emergencystretch=2em
最后考虑纤维 $f^{-1}(\pi(\infty))$，它是 $X(\Gamma)$ 中尖点（的像）的集合.其基数为 $n_\infty$，且 $\sum_{P\mapsto\pi(\infty)}e_P=d$.故它对 Hurwitz 公式的贡献为 $\sum_{P\mapsto\pi(\infty)}(e_P-1)=d-n_\infty$.
\par\endgroup



最后只需将各项合并.

\end{proof}

\begin{remark*}

关于尖点处的分歧.设 $x\in\mathbb P^1(\mathbb R)$ 是 $\Gamma$ 的一个尖点，在 $X(\Gamma)$ 中的像为 $P$.其稳定子 $\Gamma_x\{\pm1\}$ 是 $\Gamma(1)_x$ 的有限指数子群，记其指数为 $h$.用 $\Gamma(1)$ 中的元素 $\tau$ 将 $x$ 送到 $\infty$ 后，已经看到 $P$ 附近的一个局部坐标为 $\exp(2i\pi\tau(z)/h)$，而 $\pi(z)$ 附近的一个局部坐标为 $\exp(2i\pi\tau(z))$.在这些坐标下，投影 $X(\Gamma)\longrightarrow X(1)$ 因而形如 $x\mapsto x^h$，说明 $h$ 正是此投影在 $P$ 处的分歧指数.

\end{remark*}



\subsection{模解释}\label{sec:modular-1-1-8}

``模''一词指的是``模空间''，即分类某类对象的空间.在这里，曲线 $Y(1)$ 分类亏格为 $1$ 的 Riemann 曲面，分类时相差同构（即双全纯映射）的视为相同.



为解释这一点，使用（一维）\newterm{modular-complex-torus}{复环面}{complex torus}的概念.按定义，它是 $\mathbb C$ 关于格 $\Lambda=\mathbb Z\omega_1\oplus\mathbb Z\omega_2$ 的商 $\mathbb C/\Lambda$，其复结构由 $\Lambda$ 不变的全纯函数层给出，如命题~\ref{sec:modular-1-1-3} 中所述（只是这里 $\Lambda$ 的作用是正常且自由的，因而 $\mathbb C\longrightarrow\mathbb C/\Lambda$ 是局部平凡覆盖，即无分歧）.显然 $\mathbb C/\Lambda$ 紧且亏格为 $1$.反过来：

\begin{proposition*}

每个亏格为 $1$ 的 Riemann 曲面 $X$ 都双全纯同构于一个（一维）复环面.

\end{proposition*}

\begin{proof}

Riemann 单值化定理告诉我们，$X$ 的万有覆盖是 $\mathbb C$ 或 $\mathbb H$，且 $X$ 可以实现为其万有覆盖关于某个在其上正常且自由作用的群的商.



若 $X$ 形如 $\Gamma\backslash\mathbb H$，其中 $\Gamma\subset\operatorname{Aut}(\mathbb H)$，则：

\begin{itemize}

\item $\Gamma$ 必须离散（才能正常作用）；

\item $\Gamma$ 没有椭圆点（才能自由作用），也没有尖点（否则可以扩大已经紧的 $X$）；

\item $\Gamma$ 同构于 $X$ 的基本群，即 $\mathbb Z^2$.

\end{itemize}

因此，$\Gamma$ 由双曲元素组成，这些元素在 $\mathbb R^2$ 中有两条公共特征直线（同时对角化）.故共轭之后，$\Gamma$ 是 $\mathrm{SL}_2(\mathbb R)$ 中对角矩阵群（模去 $\pm1$）的子群，从而是 $\mathbb R^\times$ 的子群.但后者没有同构于 $\mathbb Z^2$ 的离散子群.矛盾.



所以 $X$ 形如 $\Gamma\backslash\mathbb C$.而 $\operatorname{Aut}(\mathbb C)$ 由映射 $z\mapsto az+b$ 组成，其中 $(a,b)\in\mathbb C^\times\times\mathbb C$.若 $a\ne1$，这样的自同构有不动点，因而不能自由作用.所以最终 $\Gamma$ 是由上述格 $\Lambda$ 给出的离散平移群，并且 $\Gamma\backslash\mathbb C=\mathbb C/\Lambda$.

\end{proof}

\begin{remark*}

一个推论是：每个亏格为 $1$ 的 Riemann 曲面都可以赋予解析群结构，因为 $\mathbb C/\Lambda$ 具有这样的结构.

\end{remark*}

现在研究两个复环面之间的双全纯映射.

\begin{proposition*}

每个全纯映射 $\varphi:\mathbb C/\Lambda\longrightarrow\mathbb C/\Lambda'$ 都形如 $[z+\Lambda]\mapsto[\alpha z+\beta+\Lambda']$，其中 $\beta\in\mathbb C$，且 $\alpha\in\mathbb C$ 满足 $\alpha\Lambda\subset\Lambda'$.

\end{proposition*}

\begin{proof}

记 $\pi:\mathbb C\longrightarrow\mathbb C/\Lambda$ 和 $\pi':\mathbb C\longrightarrow\mathbb C/\Lambda'$ 为典范投影.由于 $\mathbb C$ 单连通，存在连续映射 $\widetilde\varphi:\mathbb C\longrightarrow\mathbb C$，使 $\varphi\circ\pi=\pi'\circ\widetilde\varphi$；而且此映射全纯.对每个 $\lambda\in\Lambda$，映射 $z\mapsto\widetilde\varphi(z+\lambda)-\widetilde\varphi(z)$ 的值在离散集 $\Lambda'$ 中，由 $\mathbb C$ 连通可知它是常值.因此导数 $\widetilde\varphi'$ 是 $\Lambda$ 不变的，便下降为 $\mathbb C/\Lambda$ 上的全纯函数.该函数只能是常值，设其值为 $\alpha\in\mathbb C$.于是存在 $\beta\in\mathbb C$ 使 $\widetilde\varphi(z)=\alpha z+\beta$.

\end{proof}

\begin{remark*}

一个推论是，两个复环面之间将 $0$ 映到 $0$ 的每个全纯映射都是解析群同态.因此，若 $X$ 是亏格为 $1$ 的紧 Riemann 曲面，并给定一点 $x$，则上一注中以 $x$ 为单位元的解析群结构是唯一的，即它与所选的双全纯映射 $\mathbb C/\Lambda\xrightarrow\sim X$ 无关，只要该映射将 $0$ 映到 $x$.

\end{remark*}

格 $\Lambda$ 作为 $\mathbb Z$ 模有一组基 $(\omega_1,\omega_2)$.所有基都可以通过 $\mathrm{GL}_2(\mathbb Z)$ 在列向量 $\binom{\omega_1}{\omega_2}\in\mathbb C^2$ 上的作用相互得到.可以排列基的次序，使 $\operatorname{Im}(\omega_1/\omega_2)>0$.如此排序的基通过 $\mathrm{SL}_2(\mathbb Z)$ 的作用相互得到.由此得到格的集合 $\mathcal R$ 与集合 $\mathrm{SL}_2(\mathbb Z)\backslash\mathcal V$ 的双射，其中 $\mathcal V$ 是满足 $\omega_1/\omega_2\in\mathbb H$ 的列向量 $\binom{\omega_1}{\omega_2}\in\mathbb C^2$ 所成的集合，后一个集合表示模去 $\mathrm{SL}_2(\mathbb Z)$ 的作用.在此双射下，$\mathbb C^\times$ 在 $\mathcal R$ 上的位似作用对应于在 $\mathcal V$ 上的位似作用.刚才已经看到，商 $\mathcal R/\mathbb C^*$ 等同于复环面的同构类的集合 $\mathcal E$.另一方面，映射 $\binom{\omega_1}{\omega_2}\mapsto\omega_1/\omega_2$ 诱导 $\mathcal V/\mathbb C^\times$ 到 $\mathbb H$ 的双射.最终得到以下结果.

\begin{proposition*}

有如下双射：

\begin{center}

\begin{tikzcd}[column sep=large,row sep=large]

Y(1)=\mathrm{SL}_2(\mathbb Z)\backslash\mathbb H \arrow[d,"\sim"',"z\mapsto\mathbb Zz\oplus\mathbb Z"] & \{\text{亏格为 }1\text{ 的 Riemann 曲面}\}_{/\text{双全纯同构}}\\

\{\mathbb C\text{ 中的格}\}_{/\text{位似}} \arrow[r,"\Lambda\mapsto\mathbb C/\Lambda","\sim"'] & \{\text{复环面}\}_{/\text{同构}} \arrow[u,"\sim"']

\end{tikzcd}

\end{center}

\end{proposition*}

